\section{讨论与后续工作}
\label{sec:discussion}
\subsection{结论的边界}
单次比较不能支持普遍结论。正式论文需要重复测量，
并说明观察结果能否推广到其他条件。
\subsection{记录可复核的过程}
正式写作时，应把测量表格、计算脚本和使用条件一起保存。
修改一个数值后，检查公式、表格和正文是否同步更新；
删除章节前，检查其他位置是否仍引用该章节的标签。

本项目将三个章节保存在 chapters 文件夹中。编译入口始终是
main.tex，章节文件不需要重复声明文档类型。这样可以在单独
编辑某一章的同时，保持目录和参考文献的一致性。
\subsection{下一步检查}
\begin{enumerate}
\item 检查主文件的章节引用与编译顺序。
\item 修改表格后，核对正文和公式中的数字。
\item 大幅调整章节前保存项目备份。
\end{enumerate}
